% บทที่ 4 ผลการดำเนินงาน
% ตัวเลขทุกตัวมาจากสคริปต์ใน ../codingpy/ ที่รันกับผู้ผ่านข้อดักความใส่ใจ 107 คน
% (passed_attention ใน config.py) ยกเว้นหัวข้อสุดท้ายที่รันกับผู้ตอบทั้ง 121 คนเพื่อเทียบ
\chapter{ผลการดำเนินงาน}

\section{ข้อมูลที่ได้}

แบบสอบถามได้คำตอบ 121 ชุด ผู้ตอบ 14 คน (ร้อยละ 11.6) ไม่ผ่านข้อดักความใส่ใจและถูกตัดออก
เหลือ 107 คนที่ใช้วิเคราะห์ในบทนี้

ขั้นทำความสะอาดพบค่าที่ใช้ไม่ได้สองแบบ
ผู้ตอบ 1 คนตอบชั่วโมงการนอนเป็น 210 ซึ่งเป็นไปไม่ได้ จึงกำหนดเป็นค่าว่าง
และผู้ตอบ 3 คนตอบข้อเกรดเฉลี่ยเป็น ``-'' จึงไม่มีค่าเกรดเฉลี่ย
ตัวแปรเชิงปริมาณ 12 ตัวของผู้ตอบ 107 คนมีค่าว่างรวม 4 ช่องจาก 1,284 ช่อง
การถดถอยพหุคูณซึ่งต้องใช้แถวที่ครบทุกตัวแปรจึงเหลือ 103 คน

\section{คุณภาพและที่มาของข้อมูล}

ข้อมูลที่เก็บเองไม่มีลักษณะของข้อมูลสังเคราะห์ตามเกณฑ์ในบทที่ 3
ชั่วโมงการใช้งานร้อยละ 65.4 เป็นจำนวนเต็ม และร้อยละ 90.7 ลงท้ายด้วย .0 หรือ .5
ซึ่งเป็นแบบที่คนจริงตอบ ไม่ได้กระจายเท่ากันทุกหลัก
ข้อมูลมีค่าว่างอยู่บ้างตามที่รายงานข้างต้น
และสหสัมพันธ์ระหว่างตัวแปรต้นสูงสุดคือ 0.45 ระหว่างชั่วโมงการนอนกับคุณภาพการนอน
ซึ่งเป็นสองเรื่องที่ควรสัมพันธ์กันอยู่แล้ว ไม่มีคู่ใดสูงถึงระดับ 0.80 ถึง 0.95 แบบชุดข้อมูลสังเคราะห์

\section{ลักษณะของกลุ่มตัวอย่าง}

ผู้ตอบเป็นหญิง 59 คน (ร้อยละ 55.1) และชาย 48 คน (ร้อยละ 44.9)
อายุเฉลี่ย 19.36 ปี (SD = 0.82, ต่ำสุด 17 ปี สูงสุด 23 ปี)
ส่วนใหญ่เรียนชั้นปีที่ 2 คือ 87 คน (ร้อยละ 81.3) รองลงมาคือชั้นปีที่ 1 จำนวน 16 คน (ร้อยละ 15.0)
และชั้นปีที่ 3 จำนวน 4 คน (ร้อยละ 3.7)
ผู้ตอบ 40 คน (ร้อยละ 37.4) ไม่ได้ระบุคณะ ในกลุ่มที่ระบุ
คณะที่มีผู้ตอบมากที่สุดคือวิทยาลัยการคอมพิวเตอร์ 18 คน
รองลงมาคือคณะมนุษยศาสตร์และสังคมศาสตร์ 10 คน และคณะวิทยาศาสตร์ 8 คน

แพลตฟอร์มที่ใช้มากที่สุดคือ TikTok 45 คน (ร้อยละ 42.1) Instagram 29 คน (ร้อยละ 27.1)
และ YouTube 27 คน (ร้อยละ 25.2) ส่วน X และ Facebook มีอย่างละ 3 คน
ตารางที่~\ref{tab:content} แสดงประเภทเนื้อหาที่ดูมากที่สุด
บันเทิงมีผู้ตอบเกินครึ่ง ขณะที่ข่าวสารบ้านเมืองและการเมืองมีเพียง 8 คน

\begin{table}[H]
\centering
\caption{ประเภทเนื้อหาที่ดูมากที่สุด และกลุ่มที่ใช้วิเคราะห์}
\label{tab:content}
\begin{tabular}{|l|r|r|l|}
\hline
\textbf{ตัวเลือกในแบบสอบถาม} & \textbf{คน} & \textbf{ร้อยละ} & \textbf{กลุ่มที่ใช้วิเคราะห์} \\
\hline
บันเทิง & 58 & 54.2 & บันเทิง (58 คน) \\
\hline
ชีวิตประจำวันของคนอื่น & 13 & 12.1 & ชีวิตคนอื่น (13 คน) \\
\hline
ข่าวสารบ้านเมืองและการเมือง & 8 & 7.5 & ข่าว (8 คน) \\
\hline
เกมและอีสปอร์ต & 11 & 10.3 & \multirow{5}{*}{อื่น~ๆ (28 คน)} \\
\cline{1-3}
ความรู้และการพัฒนาตนเอง & 7 & 6.5 & \\
\cline{1-3}
กีฬา & 6 & 5.6 & \\
\cline{1-3}
สินค้าและรีวิว & 3 & 2.8 & \\
\cline{1-3}
อื่น~ๆ & 1 & 0.9 & \\
\hline
\end{tabular}
\end{table}

\section{สถิติเชิงพรรณนา}

\begin{table}[H]
\centering
\caption{สถิติเชิงพรรณนาของตัวแปรเชิงปริมาณ}
\label{tab:descriptive}
\begin{tabular}{|l|r|r|r|r|r|r|}
\hline
\textbf{ตัวแปร} & \textbf{n} & \textbf{เฉลี่ย} & \textbf{มัธยฐาน} & \textbf{SD} & \textbf{ต่ำสุด} & \textbf{สูงสุด} \\
\hline
ชั่วโมงใช้โซเชียลมีเดียต่อวัน & 107 & 7.83 & 7.20 & 2.70 & 4.0 & 15.5 \\
วันที่ติดตามข่าวเชิงลบ (ต่อ 7 วัน) & 107 & 3.46 & 3.00 & 1.54 & 1 & 7 \\
ชั่วโมงการนอนต่อคืน & 106 & 6.21 & 6.00 & 1.54 & 2 & 11 \\
คุณภาพการนอน (1--10) & 107 & 5.54 & 5.00 & 2.16 & 1 & 10 \\
วันที่ออกกำลังกายต่อสัปดาห์ & 107 & 2.00 & 2.00 & 1.81 & 0 & 7 \\
ความกดดันจากการเรียน (1--10) & 107 & 6.12 & 7.00 & 2.26 & 1 & 10 \\
เกรดเฉลี่ยสะสม & 104 & 3.20 & 3.28 & 0.49 & 2.00 & 4.00 \\
ความเพียงพอทางการเงิน (1--10) & 107 & 6.50 & 7.00 & 2.13 & 1 & 10 \\
\hline
ภาวะซึมเศร้า (0--42) & 107 & 9.42 & 6.00 & 8.13 & 0 & 36 \\
ความวิตกกังวล (0--42) & 107 & 11.40 & 10.00 & 7.47 & 0 & 38 \\
ความเครียด (0--42) & 107 & 11.89 & 10.00 & 7.78 & 0 & 34 \\
\hline
\end{tabular}
\end{table}

ผู้ตอบใช้โซเชียลมีเดียเฉลี่ย 7.83 ชั่วโมงต่อวัน ไม่มีใครใช้น้อยกว่า 4 ชั่วโมง
และนอนเฉลี่ย 6.21 ชั่วโมงต่อคืน
คะแนน DASS-21 ทั้งสามมิติเบ้ขวา ค่าเฉลี่ยสูงกว่ามัธยฐาน
แปลว่าคนส่วนใหญ่มีคะแนนต่ำถึงปานกลาง และมีคนส่วนน้อยที่คะแนนสูงมาก

\begin{table}[H]
\centering
\caption{จำนวนผู้ตอบ (ร้อยละ) ในแต่ละระดับความรุนแรงของ DASS-21}
\label{tab:severity-result}
\begin{tabular}{|l|r|r|r|}
\hline
\textbf{ระดับ} & \textbf{ซึมเศร้า} & \textbf{วิตกกังวล} & \textbf{เครียด} \\
\hline
ปกติ & 63 (58.9) & 34 (31.8) & 72 (67.3) \\
เล็กน้อย & 15 (14.0) & 16 (15.0) & 19 (17.8) \\
ปานกลาง & 16 (15.0) & 28 (26.2) & 9 (8.4) \\
รุนแรง & 10 (9.3) & 13 (12.1) & 6 (5.6) \\
รุนแรงมาก & 3 (2.8) & 16 (15.0) & 1 (0.9) \\
\hline
\end{tabular}
\end{table}

ความวิตกกังวลเป็นมิติที่ผู้ตอบอยู่นอกระดับปกติมากที่สุด คือ 73 คน (ร้อยละ 68.2)
และร้อยละ 15.0 อยู่ในระดับรุนแรงมาก
ส่วนภาวะซึมเศร้าและความเครียด ผู้ตอบส่วนใหญ่อยู่ในระดับปกติ
ตัวเลขนี้เป็นผลคัดกรองเบื้องต้นของกลุ่มตัวอย่าง ไม่ใช่การวินิจฉัย

\section{ความเชื่อมั่นของแบบวัด}
\label{sec:reliability}

\begin{table}[H]
\centering
\caption{ค่าแอลฟาของครอนบาคของ DASS-21 ($n = 107$)}
\label{tab:alpha}
\begin{tabular}{|l|c|c|c|}
\hline
\textbf{มิติ} & \textbf{จำนวนข้อ} & \textbf{แอลฟา} & \textbf{95\% CI (bootstrap)} \\
\hline
ภาวะซึมเศร้า & 7 & 0.844 & 0.790--0.880 \\
ความวิตกกังวล & 7 & 0.737 & 0.627--0.810 \\
ความเครียด & 7 & 0.771 & 0.690--0.824 \\
รวมทั้งฉบับ & 21 & 0.894 & 0.857--0.917 \\
\hline
\end{tabular}
\end{table}

ทั้งสามมิติมีแอลฟาตั้งแต่ 0.70 ขึ้นไป (ตารางที่~\ref{tab:alpha})
ภาวะซึมเศร้ามีความเชื่อมั่นสูงที่สุด ความวิตกกังวลต่ำที่สุดและขอบล่างของช่วงความเชื่อมั่นอยู่ที่ 0.63
ข้อที่สัมพันธ์กับข้ออื่นในมิติวิตกกังวลต่ำที่สุดคือข้อ 2 (ปากแห้ง, $r = 0.30$) และข้อ 9 ($r = 0.29$)
แต่ถ้าตัดข้อใดข้อหนึ่งออก แอลฟาขึ้นเพียงเป็น 0.742 จึงใช้ครบทั้ง 7 ข้อ

คะแนนสามมิติสัมพันธ์กันเองปานกลางถึงค่อนข้างสูง
ซึมเศร้ากับความเครียด $r = 0.68$ ซึมเศร้ากับวิตกกังวล $r = 0.55$ และวิตกกังวลกับความเครียด $r = 0.55$
สามมิติจึงแยกจากกันได้ แต่ซึมเศร้ากับความเครียดซ้อนกันพอสมควร

\section{ผลตามคำถามวิจัย}

\subsection{คำถามข้อ 1 และ 2: เวลาที่ใช้กับคะแนนทั้งสามมิติ}

ชั่วโมงใช้โซเชียลมีเดียต่อวันไม่สัมพันธ์กับมิติใดอย่างมีนัยสำคัญ (ตารางที่~\ref{tab:q123})
ค่า $r$ ทั้งสามค่าติดลบเล็กน้อย คือคนที่ใช้มากกว่ามีคะแนนต่ำกว่าเล็กน้อย
แต่ช่วงความเชื่อมั่น 95\% ทุกช่วงคร่อมศูนย์
และ $r^2$ สูงสุดเพียง 0.022 หรือเวลาที่ใช้อธิบายความแปรปรวนของคะแนนได้ไม่ถึงร้อยละ 2.2
สหสัมพันธ์สเปียร์แมนให้ผลทางเดียวกัน ($\rho$ อยู่ระหว่าง $-0.12$ ถึง $-0.03$)
และเมื่อควบคุมตัวแปรอื่นแล้วก็ยังไม่มีนัยสำคัญในมิติใด

คำถามข้อ 2 ถามว่าโซเชียลมีเดียสัมพันธ์กับมิติใดมากที่สุด
ค่า $|r|$ สูงสุดอยู่ที่ความวิตกกังวล แต่ทั้งสามมิติมีขนาดใกล้กันและไม่มีมิติใดมีนัยสำคัญ
จึงยังตอบไม่ได้ว่ามิติใดสัมพันธ์มากกว่ากัน

\begin{table}[H]
\centering
\caption{สหสัมพันธ์ระหว่างพฤติกรรมการใช้โซเชียลมีเดียกับคะแนน DASS-21 ($n = 107$)}
\label{tab:q123}
\begin{tabular}{|l|l|c|c|c|c|}
\hline
\textbf{ตัวแปรต้น} & \textbf{มิติ} & \textbf{$r$} & \textbf{95\% CI} & \textbf{$p$} & \textbf{partial $r$ ($p$)} \\
\hline
\multirow{3}{*}{ชั่วโมงใช้งานต่อวัน} & ซึมเศร้า & $-0.136$ & $-0.293$, $0.033$ & 0.161 & $-0.101$ (0.322) \\
 & วิตกกังวล & $-0.148$ & $-0.334$, $0.030$ & 0.129 & $-0.130$ (0.200) \\
 & เครียด & $-0.087$ & $-0.258$, $0.094$ & 0.376 & $-0.064$ (0.532) \\
\hline
\multirow{3}{*}{วันที่ติดตามข่าว} & ซึมเศร้า & $-0.013$ & $-0.187$, $0.168$ & 0.892 & $-0.001$ (0.993) \\
 & วิตกกังวล & $0.146$ & $-0.034$, $0.328$ & 0.134 & $0.166$ (0.103) \\
 & เครียด & $0.101$ & $-0.074$, $0.277$ & 0.303 & $0.106$ (0.297) \\
\hline
\end{tabular}
\end{table}

\noindent
partial $r$ ควบคุมชั่วโมงการนอน คุณภาพการนอน ความกดดันจากการเรียน เกรดเฉลี่ย และความเพียงพอทางการเงิน
($n = 103$)

\begin{figure}[H]
\centering
\includegraphics[width=\textwidth]{qr2-social-media-dass-dimension.png}
\caption{สหสัมพันธ์ของเวลาใช้งานและการติดตามข่าวกับ DASS-21 (เส้นคือ 95\% CI)}
\label{fig:qr2}
\end{figure}

\subsection{คำถามข้อ 3: การติดตามข่าวเชิงลบกับความวิตกกังวล}

จำนวนวันที่ติดตามข่าวเชิงลบสัมพันธ์ทางบวกกับความวิตกกังวลมากกว่ามิติอื่น
($r = 0.146$) ซึ่งเป็นทิศทางที่คาดไว้
แต่ไม่มีนัยสำคัญ ($p = 0.134$) และเมื่อควบคุมตัวแปรอื่นแล้วก็ยังไม่มีนัยสำคัญ
(partial $r = 0.166$, $p = 0.103$) ส่วนภาวะซึมเศร้าแทบไม่สัมพันธ์เลย ($r = -0.013$)
รูปแบบของผลจึงเข้ากับสมมติฐาน แต่ข้อมูลยังไม่มากพอจะยืนยัน

\subsection{คำถามข้อ 4: ประเภทเนื้อหาที่เสพ}

คะแนนของผู้ที่ดูเนื้อหาต่างกลุ่มกันไม่ต่างกันในทุกมิติ
(Kruskal-Wallis: ซึมเศร้า $H = 0.418$, $p = 0.936$;
วิตกกังวล $H = 0.659$, $p = 0.883$; เครียด $H = 1.460$, $p = 0.691$)
กลุ่มข่าวมีคะแนนความวิตกกังวลเฉลี่ยสูงที่สุด (13.00) เทียบกับบันเทิง 11.59 อื่น~ๆ 11.00
และชีวิตคนอื่น 10.46 แต่ ANOVA ไม่พบความต่าง ($F(3, 103) = 0.225$, $p = 0.879$, $\eta^2 = 0.007$)
กลุ่มข่าวมีเพียง 8 คน จึงมีกำลังในการตรวจความต่างเล็ก~ๆ ต่ำมาก

\begin{figure}[H]
\centering
\includegraphics[width=\textwidth]{q4-content-type.png}
\caption{คะแนน DASS-21 แยกตามกลุ่มเนื้อหาที่ดูมากที่สุด}
\label{fig:content}
\end{figure}

\subsection{คำถามข้อ 5: เวลาที่ใช้กับการนอน}

ชั่วโมงใช้โซเชียลมีเดียไม่สัมพันธ์กับชั่วโมงการนอน ($r = 0.049$, $p = 0.616$, $n = 106$)
และไม่สัมพันธ์กับคุณภาพการนอน ($r = 0.114$, $p = 0.243$)
ในกลุ่มตัวอย่างนี้ คนที่ใช้มากจึงไม่ได้นอนน้อยกว่าหรือนอนแย่กว่า

แต่คุณภาพการนอนเองสัมพันธ์กับคะแนน DASS-21 (รูปที่~\ref{fig:lifestyle})
คนที่ให้คะแนนคุณภาพการนอนสูงมีภาวะซึมเศร้าต่ำกว่า ($r = -0.276$, $p = 0.004$)
และความเครียดต่ำกว่า ($r = -0.320$, $p < 0.001$)
ส่วนชั่วโมงการนอนและจำนวนวันที่ออกกำลังกายไม่สัมพันธ์กับมิติใดอย่างมีนัยสำคัญ

\begin{figure}[H]
\centering
\includegraphics[width=0.8\textwidth]{q5-interaction-dass.png}
\caption{สหสัมพันธ์ระหว่างพฤติกรรมในชีวิตประจำวันกับคะแนน DASS-21 แต่ละมิติ}
\label{fig:lifestyle}
\end{figure}

\subsection{คำถามข้อ 6: เมื่อควบคุมปัจจัยอื่น}

ตัวแปรต้นทั้งแปดตัวไม่มีปัญหาสัมพันธ์กันเอง ค่า VIF อยู่ระหว่าง 1.05 ถึง 1.54 ทุกตัวต่ำกว่า 5
ความกดดันจากการเรียนสัมพันธ์กับคะแนนความเครียด ($r = 0.280$, $p = 0.004$)
แต่ไม่สูงถึงขั้นที่ต้องตัดออกตามแผนในเค้าโครง

\begin{table}[H]
\centering
\caption{การถดถอยแบบลำดับขั้น: $R^2$ ที่เพิ่มขึ้นเมื่อใส่ตัวแปรโซเชียลมีเดีย ($n = 103$)}
\label{tab:hierarchical}
\begin{tabular}{|l|c|c|c|c|c|c|}
\hline
\textbf{มิติ} & \textbf{$R^2$ ขั้น 1} & \textbf{$R^2$ ขั้น 2} & \textbf{$\Delta R^2$} & \textbf{F-change} & \textbf{$p$} & \textbf{adj. $R^2$} \\
\hline
ซึมเศร้า & 0.263 & 0.271 & 0.008 & 0.489 & 0.615 & 0.208 \\
วิตกกังวล & 0.058 & 0.103 & 0.045 & 2.355 & 0.100 & 0.027 \\
เครียด & 0.191 & 0.205 & 0.013 & 0.796 & 0.454 & 0.137 \\
\hline
\end{tabular}
\end{table}

ขั้นที่ 1 ใส่เฉพาะตัวแปรควบคุมหกตัว ขั้นที่ 2 เพิ่มชั่วโมงใช้งานและวันที่ติดตามข่าว
ตารางที่~\ref{tab:hierarchical} แสดงว่าตัวแปรโซเชียลมีเดียทั้งสองตัวเพิ่ม $R^2$ ได้น้อย
และไม่มีนัยสำคัญในมิติใด
มิติวิตกกังวลเพิ่มมากที่สุด ($\Delta R^2 = 0.045$, $p = 0.100$)
แต่หลังปรับ Bonferroni ค่า $p$ เป็น 0.301
คำตอบของคำถามข้อ 6 คือเมื่อควบคุมปัจจัยอื่นแล้ว
ไม่พบว่าเวลาที่ใช้กับโซเชียลมีเดียอธิบายคะแนนเพิ่มเติม

\begin{table}[H]
\centering
\caption{สัมประสิทธิ์ $b$ (ค่า $p$) ของสมการเต็ม SE แบบ HC3 ($n = 103$)}
\label{tab:coef}
\begin{tabular}{|l|r|r|r|}
\hline
\textbf{ตัวแปรต้น} & \textbf{ซึมเศร้า} & \textbf{วิตกกังวล} & \textbf{เครียด} \\
\hline
ชั่วโมงใช้งานต่อวัน & $-0.268$ (0.299) & $-0.391$ (0.153) & $-0.187$ (0.485) \\
วันที่ติดตามข่าว & $0.023$ (0.961) & $0.872$ (0.079) & $0.534$ (0.297) \\
ชั่วโมงการนอน & $1.150$ (0.065) & $0.069$ (0.901) & $0.195$ (0.725) \\
คุณภาพการนอน & $-0.921$ (0.027) & $-0.162$ (0.663) & $-0.878$ (0.044) \\
วันที่ออกกำลังกาย & $0.176$ (0.692) & $-0.066$ (0.882) & $-0.185$ (0.669) \\
ความกดดันจากการเรียน & $0.979$ (0.009) & $0.329$ (0.394) & $0.800$ (0.037) \\
เกรดเฉลี่ยสะสม & $-0.706$ (0.655) & $-0.958$ (0.528) & $0.115$ (0.942) \\
ความเพียงพอทางการเงิน & $-1.125$ (0.001) & $-0.663$ (0.048) & $-0.626$ (0.079) \\
\hline
$R^2$ & 0.271 & 0.103 & 0.205 \\
\hline
\end{tabular}
\end{table}

ตารางที่~\ref{tab:coef} แสดงว่าตัวแปรที่มีน้ำหนักในสมการเป็นตัวแปรควบคุม ไม่ใช่ตัวแปรโซเชียลมีเดีย
ในสมการภาวะซึมเศร้า ความเพียงพอทางการเงินมีน้ำหนักมากที่สุด
(เงินพอขึ้น 1 ระดับ คะแนนลดลง 1.125 คะแนน, $\beta = -0.293$)
ตามด้วยความกดดันจากการเรียน ($\beta = 0.269$) และคุณภาพการนอน ($\beta = -0.242$)
ในสมการความเครียด คุณภาพการนอน ($\beta = -0.243$) และความกดดันจากการเรียน ($\beta = 0.232$) มีน้ำหนักมากที่สุด
เมื่อปรับ Bonferroni แล้ว เหลือเพียงความเพียงพอทางการเงิน ($p = 0.002$)
และความกดดันจากการเรียน ($p = 0.026$) ในสมการภาวะซึมเศร้าที่ยังมีนัยสำคัญ
สมการความวิตกกังวลอธิบายได้น้อยที่สุด (adj. $R^2 = 0.027$)

การตรวจข้อตกลงเบื้องต้น (รูปที่~\ref{fig:diagnostics}) พบว่าค่าคลาดเคลื่อนของสมการภาวะซึมเศร้า
มีความแปรปรวนไม่คงที่ (Breusch-Pagan $p = 0.013$) ซึ่งเป็นเหตุผลที่ใช้ค่าความคลาดเคลื่อนมาตรฐานแบบ HC3
และค่าคลาดเคลื่อนของสมการความวิตกกังวลไม่แจกแจงปกติ (Shapiro-Wilk $p = 0.001$)
ปลายด้านขวาของ Q-Q plot เชิดขึ้นจากผู้ตอบไม่กี่คนที่คะแนนสูงมาก
ส่วนสมการอื่นไม่พบปัญหาทั้งสองข้อ ($p > 0.15$ ทุกค่า)

\begin{figure}[H]
\centering
\includegraphics[width=\textwidth]{q0.png}
\caption{การตรวจข้อตกลงเบื้องต้นของสมการถดถอย แถวบนคือค่าคลาดเคลื่อนเทียบกับค่าที่โมเดลทำนาย แถวล่างคือ Q-Q plot}
\label{fig:diagnostics}
\end{figure}

\subsection{คำถามข้อ 7: ความเครียดเมื่อผลการเรียนและการเงินเท่ากัน}

ในสมการความเครียด ซึ่งควบคุมเกรดเฉลี่ย ความเพียงพอทางการเงิน และตัวแปรอื่นทั้งหมดแล้ว
สัมประสิทธิ์ของชั่วโมงใช้งานคือ $b = -0.187$ คะแนนต่อชั่วโมง
(95\% CI $-0.713$ ถึง $0.338$, $p = 0.485$, หลังปรับ Bonferroni $p = 1.000$)
ในกลุ่มนักศึกษาที่ผลการเรียนและการเงินใกล้เคียงกัน
คนที่ใช้โซเชียลมีเดียมากกว่าจึงไม่ได้มีคะแนนความเครียดต่างจากคนที่ใช้น้อยกว่า
ช่วงความเชื่อมั่นบอกว่าความต่างจริงน่าจะอยู่ระหว่าง $-0.71$ ถึง $+0.34$
คะแนนต่อชั่วโมง จากคะแนนเต็ม 42

\begin{figure}[H]
\centering
\includegraphics[width=0.85\textwidth]{q7-partial-regression-stress.png}
\caption{Partial regression plot ของชั่วโมงใช้งานกับคะแนนความเครียด หลังหักผลของตัวแปรต้นตัวอื่น ความชันของเส้นเท่ากับค่า $b$ ในตารางที่~\ref{tab:coef}}
\label{fig:partial}
\end{figure}

\section{สภาพการเงินและผลการเรียน}

การวิเคราะห์เพิ่มเติมนี้แบ่งผู้ตอบตามคะแนนความเพียงพอทางการเงินเป็น 3 กลุ่ม
คือเงินไม่พอใช้ (1--4, 18 คน) พอใช้ (5--7, 54 คน) และสบาย (8--10, 35 คน)
แล้วเทียบคะแนนเฉลี่ยด้วย ANOVA (ตารางที่~\ref{tab:finance} และรูปที่~\ref{fig:finance})

\begin{table}[H]
\centering
\caption{คะแนนเฉลี่ย DASS-21 แยกตามสภาพการเงิน ($n = 107$)}
\label{tab:finance}
\begin{tabular}{|l|c|c|c|c|c|c|}
\hline
\textbf{มิติ} & \textbf{ไม่พอใช้} & \textbf{พอใช้} & \textbf{สบาย} & \textbf{$F(2, 104)$} & \textbf{$p$} & \textbf{$\eta^2$} \\
\hline
ซึมเศร้า & 16.6 & 8.9 & 6.6 & 10.91 & $< 0.001$ & 0.173 \\
วิตกกังวล & 14.3 & 11.2 & 10.2 & 1.91 & 0.153 & 0.035 \\
เครียด & 16.9 & 11.3 & 10.2 & 5.02 & 0.008 & 0.088 \\
\hline
\end{tabular}
\end{table}

ภาวะซึมเศร้าต่างกันตามสภาพการเงินมากที่สุด
กลุ่มเงินไม่พอใช้มีคะแนนเฉลี่ยสูงกว่ากลุ่มสบาย 10.0 คะแนน (Cohen's $d = 1.59$)
และสภาพการเงินอธิบายความแปรปรวนของคะแนนซึมเศร้าได้ร้อยละ 17.3
ความเครียดต่างกันเช่นกัน ($d = 0.95$) ส่วนความวิตกกังวลไม่ต่างกันอย่างมีนัยสำคัญ
การทดสอบ Kruskal-Wallis ให้ข้อสรุปเดียวกันทั้งสามมิติ
เมื่อใช้คะแนนความเพียงพอ 1--10 โดยตรงแทนการแบ่งกลุ่ม
ก็สัมพันธ์ทางลบกับภาวะซึมเศร้า ($r = -0.374$) แรงกว่าชั่วโมงใช้โซเชียลมีเดีย ($r = -0.136$) เกือบสามเท่า

\begin{figure}[H]
\centering
\includegraphics[width=\textwidth]{qr7-finance-dass-dimension.png}
\caption{คะแนนเฉลี่ย DASS-21 แยกตามสภาพการเงิน เส้นบนแท่งคือค่าคลาดเคลื่อนมาตรฐาน}
\label{fig:finance}
\end{figure}

เกรดเฉลี่ยสะสมไม่สัมพันธ์กับมิติใดเลย ($|r| \le 0.063$, $p > 0.5$, $n = 104$)

\section{การทดสอบสมมติฐานตามที่รายวิชากำหนด}

\begin{table}[H]
\centering
\caption{ผลการทดสอบสมมติฐานสี่แบบ ($\alpha = 0.05$)}
\label{tab:hypothesis}
\begin{tabularx}{\textwidth}{|L|l|l|l|}
\hline
\textbf{สมมติฐาน} & \textbf{การทดสอบ} & \textbf{ผล} & \textbf{ข้อสรุป} \\
\hline
ชั่วโมงการนอนเฉลี่ยเท่ากับ 7 ชั่วโมง & one-sample $t$ & $t(105) = -5.244$, $p < 0.001$ & ปฏิเสธ $H_0$ \\
\hline
สัดส่วนผู้ที่ความวิตกกังวลเกินระดับปกติเท่ากับ 0.5 & one-proportion $z$ & $z = 3.770$, $p < 0.001$ & ปฏิเสธ $H_0$ \\
\hline
คะแนนซึมเศร้าเฉลี่ยของชายและหญิงเท่ากัน & independent $t$ & $t(105) = 1.048$, $p = 0.297$ & ไม่ปฏิเสธ $H_0$ \\
\hline
คะแนนวิตกกังวลเฉลี่ยของ 4 กลุ่มเนื้อหาเท่ากัน & one-way ANOVA & $F(3, 103) = 0.225$, $p = 0.879$ & ไม่ปฏิเสธ $H_0$ \\
\hline
\end{tabularx}
\end{table}

ผู้ตอบนอนเฉลี่ย 6.21 ชั่วโมง (SD = 1.54, $n = 106$) น้อยกว่า 7 ชั่วโมงอย่างมีนัยสำคัญ
ข้อมูลชั่วโมงการนอนไม่แจกแจงปกติ (Shapiro-Wilk $p = 0.002$)
แต่ที่ $n = 106$ การทดสอบทียังใช้ได้ และ Wilcoxon signed-rank ให้ข้อสรุปเดียวกัน
สัดส่วนผู้ที่ความวิตกกังวลเกินระดับปกติคือ 0.682 (73 จาก 107 คน) สูงกว่าครึ่งอย่างมีนัยสำคัญ
ผู้ชายมีคะแนนซึมเศร้าเฉลี่ย 10.33 ($n = 48$) ผู้หญิง 8.68 ($n = 59$)
ต่างกัน 1.66 คะแนนแต่ไม่มีนัยสำคัญ (Cohen's $d = 0.20$, Mann-Whitney $p = 0.219$)

\section{ผลเมื่อรวมผู้ที่ไม่ผ่านข้อดักความใส่ใจ}
\label{sec:sensitivity}

เพื่อดูว่าการตัดผู้ตอบ 14 คนเปลี่ยนข้อสรุปหรือไม่ จึงรันการวิเคราะห์ชุดเดิมกับผู้ตอบทั้ง 121 คน
ข้อสรุปส่วนใหญ่เหมือนเดิม
ชั่วโมงใช้งานยังไม่สัมพันธ์กับมิติใด ($r$ อยู่ระหว่าง $-0.10$ ถึง $-0.03$)
และสภาพการเงินยังเป็นตัวแปรที่แยกคะแนนซึมเศร้าได้ชัดที่สุด ($F(2, 118) = 15.87$, $p < 0.001$)

ผลที่เปลี่ยนมีสองข้อ ข้อแรกคือการติดตามข่าวกับความวิตกกังวล
เมื่อรวม 121 คน สหสัมพันธ์บางส่วนมีนัยสำคัญ (partial $r = 0.189$, $p = 0.047$)
แต่เมื่อใช้ 107 คนไม่มีนัยสำคัญ (partial $r = 0.166$, $p = 0.103$)
ข้อที่สองคือคะแนนซึมเศร้าของชายและหญิง
เมื่อรวม 121 คนต่างกันอย่างมีนัยสำคัญ ($t = 1.992$, $p = 0.049$)
แต่เมื่อใช้ 107 คนไม่ต่าง ($p = 0.297$)
ผลทั้งสองข้อจึงขึ้นกับผู้ตอบ 14 คนที่ไม่ได้อ่านแบบสอบถามจริง และไม่ควรนำไปสรุป
